\chapter{Complexity — P Advice Subset P Poly}
%%% generated by the blueprint pipeline from TCSlib.Complexity.CircuitComplexity.PAdviceSubsetPPoly
\section{Overview}

Declarations of the Lean module \texttt{TCSlib.Complexity.CircuitComplexity.PAdviceSubsetPPoly}.

\section{Declarations}

\begin{definition}[Encoded pair layout]
\lean{Complexity.pairLayout}
\statementsource{Yap83}{
  pdf-1a201b302069-p002-b002,
  pdf-1a201b302069-p003-b002
}
\label{Complexity.pairLayout}
\leanok
\uses{BoolCircuit.BitSrc}
For natural numbers $n$ and $a$, form a list of bit sources by taking, for each
$i=0,\ldots,n-1$, two consecutive copies of input bit $i$, then the constants $0$ and
$1$, and then, for each $j=0,\ldots,a-1$, input bit $n+j$.
\end{definition}

\begin{theorem}[Length of the pair layout]
\lean{Complexity.length_pairLayout}
\label{Complexity.length_pairLayout}
\leanok
\uses{Complexity.CfgTableau.length_flatMap_const, Complexity.pairLayout}
\difficulty{3}
For natural numbers $n$ and $a$, form a list of bit sources by listing each of the first
$n$ input bits twice, then the constants $0$ and $1$, then the next $a$ input bits. The
resulting list has length $2n+2+a$.
\end{theorem}

\begin{theorem}[Validity of the pair layout]
\lean{Complexity.pairLayout_valid}
\label{Complexity.pairLayout_valid}
\leanok
\uses{BoolCircuit.BitSrc.Valid, Complexity.pairLayout}
\difficulty{4}
For any natural numbers $n$, $a$, and $w$, every bit source occurring in the pair layout
for lengths $n$ and $a$ is valid with $n+a$ input bits, block width $w$, and instruction
index $0$.
\end{theorem}

\begin{theorem}[Range map equals list prefix]
\lean{Complexity.range_map_getD}
\label{Complexity.range_map_getD}
\leanok
\difficulty{3}
Let $z$ be a finite Boolean word and let $n\in\mathbb{N}$ satisfy $n\le |z|$. The list
obtained by mapping each index $i=0,\ldots,n-1$ to the bit of $z$ at index $i$ (using
false as the default value for an out-of-range index) is the prefix of $z$ of length
$n$.
\end{theorem}

\begin{theorem}[Evaluation of the pair layout]
\lean{Complexity.pairLayout_eval}
\label{Complexity.pairLayout_eval}
\leanok
\uses{BoolCircuit.BitSrc.eval, Complexity.pairLayout, Complexity.range_map_getD, Turing.pairEncode}
\difficulty{5}
For any natural numbers $n$ and $a$, and any Boolean-valued functions $x$ on the first
$n$ indices and $y$ on the first $a$ indices, evaluating the pair layout on the
concatenation of these two finite bit strings gives the self-delimiting encoding of
those strings: the resulting bit string is $\langle x,y\rangle$.
\end{theorem}

\begin{theorem}[Reconstruction of a Boolean list from indexed values]
\lean{Complexity.ofFn_getD}
\label{Complexity.ofFn_getD}
\leanok
\difficulty{3}
Let $a\in\mathbb{N}$ and let $l$ be a Boolean list of length $a$. The list whose $j$th
entry, for $j\in\{0,\ldots,a-1\}$, is the $j$th entry of $l$ when present and false
otherwise is exactly $l$.
\end{theorem}

\begin{theorem}[Advice-size polynomial bound]
\lean{Complexity.advice_size_le}
\label{Complexity.advice_size_le}
\leanok
\difficulty{6}
For all natural numbers $k,K,c,c_0,C,d,n$, we have
\[
n+C(n+1)^d+2+(c_0(n^c+1)+1)\bigl(k(2(c_0(n^c+1))+1)+(2n+2+C(n+1)^d)+3\bigr)K\leq
\bigl(C+3+(2c_0+1)(k(4c_0+1)+C+7)K\bigr)(n+1)^{2(c+d+1)}.
\]
\end{theorem}

\begin{theorem}[Advice time yields polynomial-size circuits]
\lean{Complexity.mem_PPoly_of_mem_DTIMEAdvice}
\label{Complexity.mem_PPoly_of_mem_DTIMEAdvice}
\leanok
\uses{BoolCircuit.DAGCircuit, BoolCircuit.DAGCircuit.eval, BoolCircuit.PPoly, Complexity.DTIMEAdvice, Complexity.advice_size_le, Complexity.cfgCircuit, Complexity.cfgCircuit_eval, Complexity.cfgCircuit_isFaninTwo, Complexity.cfgCircuit_size_le, Complexity.cfgConst, Complexity.decide_emits_eq_of_computesInTime, Complexity.length_pairLayout, Complexity.ofFn_getD, Complexity.pairLayout, Complexity.pairLayout_eval, Complexity.pairLayout_valid, Language.InPPoly.of_hardwire, Turing.MultiTapeTM.indicator}
\difficulty{6}
Let $c,C,d$ be natural numbers and let $a:\mathbb N\to\mathbb N$ satisfy $a(n)\le
C(n+1)^d$ for every $n$. If a language $L$ has advice of length exactly $a(n)$ and is
decidable, on every input paired with its advice, by a finite Turing machine within a
constant multiple of $n^c+1$ steps, then $L$ has a family of Boolean circuits with
fan-in at most two and size bounded by a polynomial in the input length.
\end{theorem}

\begin{theorem}[Polynomial Advice Yields Polynomial-Size Circuits]
\lean{Complexity.PAdvicePoly_subset_PPoly}
\statementsource{Yap83}{
  pdf-1a201b302069-p003-b002,
  pdf-1a201b302069-p003-b006
}
\label{Complexity.PAdvicePoly_subset_PPoly}
\leanok
\uses{BoolCircuit.PPoly, Complexity.PAdvicePoly, Complexity.mem_PPoly_of_mem_DTIMEAdvice}
\difficulty{3}
Every binary language decidable by a finite deterministic Turing machine in polynomial
time when given an advice string of length exactly $C(n+1)^d$ for inputs of length $n$,
for some natural numbers $C,d$, has a family of Boolean circuits of size at most
$a(n+1)^k$ for some natural numbers $a,k$. Each circuit is a well-formed acyclic
directed graph of AND, OR, and NOT gates, with each NOT gate having exactly one input
and each gate reading at most two earlier vertices.
\end{theorem}

\begin{theorem}[Polynomial circuit size equals polynomial advice time]
\lean{Complexity.PPoly_eq_PAdvicePoly}
\label{Complexity.PPoly_eq_PAdvicePoly}
\leanok
\uses{Complexity.PAdvicePoly, Complexity.PAdvicePoly_subset_PPoly, Complexity.PPoly_subset_PAdvicePoly, Language.InPPoly}
\difficulty{1}
The class of Boolean languages decidable by a family of fan-in-two circuits whose sizes
are bounded by a polynomial in the input length equals the class decidable in polynomial
time by a Turing machine given an advice string depending only on the input length,
where the advice length is polynomially bounded.
\end{theorem}

\begin{theorem}[Polynomial circuit size as a union of advice classes]
\lean{Complexity.PPoly_eq_iUnion_DTIMEAdvice_pow}
\label{Complexity.PPoly_eq_iUnion_DTIMEAdvice_pow}
\leanok
\uses{Complexity.DTIMEAdvice, Complexity.mem_DTIMEAdvice_pow_of_inSIZE, Complexity.mem_PPoly_of_mem_DTIMEAdvice, Language.InPPoly}
\difficulty{4}
The set of Boolean languages having a family of fan-in-two circuits of size at most
$a(n+1)^k$ for some natural numbers $a,k$ equals the union, over natural numbers $c,d$,
of the classes decidable by a polynomial-time Turing machine with advice strings of
length exactly $n^d$ on inputs of length $n$.
\end{theorem}

\begin{theorem}[Circuit polynomial time as a union of advice classes]
\lean{Complexity.PPoly_eq_iUnion_DTIMEAdvice_pow'}
\label{Complexity.PPoly_eq_iUnion_DTIMEAdvice_pow'}
\leanok
\uses{BoolCircuit.PPoly, Complexity.DTIMEAdvice, Complexity.PPoly_eq_iUnion_DTIMEAdvice_pow}
\difficulty{1}
The class of Boolean languages decidable by polynomial-size fan-in-two circuit families
equals the union, over natural numbers $c,d$, of the classes decidable by a
polynomial-time Turing machine with advice strings of length exactly $n^d$ on inputs of
length $n$.
\end{theorem}

\begin{theorem}[Unary halting with one-bit advice]
\lean{Complexity.UHALT_mem_DTIMEAdvice_one}
\label{Complexity.UHALT_mem_DTIMEAdvice_one}
\leanok
\uses{Complexity.DTIMEAdvice, Complexity.UHALT, Complexity.UHALT_le_allOnes, Complexity.mem_DTIMEAdvice_one_of_le_allOnes, Turing.MachineCode}
\difficulty{1}
For every machine representation scheme $c$, the unary halting language associated with
$c$ is decidable in time $c_0(n+1)$ with advice of length exactly one bit, for some
constant $c_0\in\mathbb{N}$. Here the advice may depend on the input length, and the
decision machine receives the paired encoding of the input and that advice.
\end{theorem}

\begin{theorem}[Universal halting with polynomial advice]
\lean{Complexity.UHALT_mem_PAdvicePoly}
\label{Complexity.UHALT_mem_PAdvicePoly}
\leanok
\uses{Complexity.PAdvicePoly, Complexity.UHALT, Complexity.UHALT_le_allOnes, Complexity.mem_PAdvicePoly_of_le_allOnes, Turing.MachineCode}
\difficulty{1}
For every machine representation scheme $c$, the unary language consisting of numbers
whose binary expansions begin with $1s$, where the suffix $s$ encodes a pair of an
advice string and an input on which the machine decoded from that advice halts, belongs
to polynomial-time with polynomial-length advice. Thus there are fixed polynomial
bounds, an advice function of the specified polynomial length, and a finite binary
Turing machine that decides this language within the specified time bound using that
advice.
\end{theorem}
